%-------------------------------------------------------------------------------
\section{Ocean model} \label{sec:basic_usel_ocean}
%-------------------------------------------------------------------------------

The ocean component is activated and configured in \namelist{PARAM_OCEAN}.
The computational mesh for the ocean component is selected by \nmitem{OCEAN_MESH_TYPE}, 
and the ocean dynamics is enabled by \nmitem{OCEAN_DYN_DO}.

\editboxtwo{
 \verb|&PARAM_OCEAN | & \\
 \verb|ACTIVATE_FLAG   = .true.,|     & ; Enable the ocean component \\
 \verb|OCEAN_MESH_TYPE = "GLOBAL",|   & ; Type of ocean mesh: REGIONAL or GLOBAL \\
 \verb|OCEAN_DYN_DO    = .true.,|     & ; Enable ocean dynamics \\
 \verb|TIME_DT          = 3600.D0,|   & ; Time step of the ocean component \\
 \verb|TIME_DT_UNIT     = "SEC",|      & ; Unit for \verb|TIME_DT| \\
 \verb|/ | & \\
}

\nmitem{OCEAN_MESH_TYPE} can be set to \verb|"REGIONAL"| or \verb|"GLOBAL"|.


%-------------------------------------------------------------------------------
\subsection{Setting grids for ocean model}
\label{subsec:ocean_mesh}
%-------------------------------------------------------------------------------

The computational grid of the ocean model is configured in \namelist{PARAM_OCEAN_MESH}.
The available parameters depend on the mesh type selected by \nmitem{OCEAN_MESH_TYPE} in \namelist{PARAM_OCEAN}.

At present, the horizontal mesh configuration of the ocean component must be identical to that of the atmospheric component.
Therefore, the horizontal domain size, MPI decomposition, number of horizontal elements, and horizontal polynomial order should be specified consistently for both components.
This restriction arises from the present atmosphere--ocean coupling,
which exchanges two-dimensional fields without horizontal remapping.


%------------------------------------------------------------------------
\subsubsection{Regional mesh}
%------------------------------------------------------------------------

For the regional ocean model using the Cartesian coordinates,
the computational domain and DG mesh are specified in \namelist{PARAM_OCEAN_MESH}.

\editboxtwo{
 \verb|&PARAM_OCEAN_MESH | & \\
 \verb|NprcX       = 1,|       & ; Number of MPI processes in the x-direction \\
 \verb|NprcY       = 1,|       & ; Number of MPI processes in the y-direction \\
 \verb|NeX         = 2,|       & ; Number of elements in x per MPI process \\
 \verb|NeY         = 2,|       & ; Number of elements in y per MPI process \\
 \verb|NeZ         = 1,|       & ; Number of vertical elements \\
 \verb|dom_xmin    = 0.D0,|    & ; Minimum x-coordinate [m] \\
 \verb|dom_xmax    = 100.D3,|  & ; Maximum x-coordinate [m] \\
 \verb|dom_ymin    = 0.D0,|    & ; Minimum y-coordinate [m] \\
 \verb|dom_ymax    = 100.D3,|  & ; Maximum y-coordinate [m] \\
 \verb|dom_zmin    = 0.D0,|    & ; Minimum vertical coordinate [m] \\
 \verb|dom_zmax    = 50.D0,|   & ; Maximum vertical coordinate [m] \\
 \verb|PolyOrder_h = 2,|       & ; Horizontal polynomial order \\
 \verb|PolyOrder_v = 1,|       & ; Vertical polynomial order \\
 \verb|/ | & \\
}

An arbitrary non-uniform element spacing can also be specified in the vertical direction by giving the locations of all element boundaries using \nmitem{FZ}.
\editboxtwo{
 \verb|&PARAM_OCEAN_MESH | & \\
 \verb|FZ = 0.D0, 10.D0, 50.D0,| & ; Vertical element boundaries [m] \\
 \verb|/ | & \\
}

%------------------------------------------------------------------------
\subsubsection{Global mesh}
%------------------------------------------------------------------------

For the global ocean model using the cubed-sphere coordinates,
the grid is also configured in \namelist{PARAM_OCEAN_MESH}.

\editboxtwo{
 \verb|&PARAM_OCEAN_MESH | & \\
 \verb|Nprc                    = 6,|      & ; Total number of MPI processes \\
 \verb|NeGX                    = 6,|      & ; Number of elements in each cubed-sphere panel \\
 \verb|NeGY                    = 6,|      & ; Number of elements in each cubed-sphere panel \\
 \verb|NeZ                     = 1,|      & ; Number of vertical elements \\
 \verb|PolyOrder_h             = 2,|      & ; Horizontal polynomial order \\
 \verb|PolyOrder_v             = 1,|      & ; Vertical polynomial order \\
 \verb|dom_zmin                = 0.D0,|   & ; Minimum vertical coordinate [m] \\
 \verb|dom_zmax                = 50.D0,|  & ; Maximum vertical coordinate [m] \\
 \verb|SHALLOW_OCN_APPROX_FLAG = .true.,| & ; Apply shallow-ocean approximation \\
 \verb|/ | & \\
}

%-------------------------------------------------------------------------------
\subsection{Ocean-model configurations}
\label{subsec:ocean_model_config}
%-------------------------------------------------------------------------------

Currently, SCALE-DG provides two ocean-model options: the slab ocean model and a fixed sea-surface temperature option.
The slab ocean model predicts the mixed-layer temperature by solving the heat budget of a prescribed mixed layer, 
while the fixed-temperature option keeps the initial sea-surface temperature unchanged throughout the simulation.
These are configured in \namelist{PARAM_OCEAN_DYN}.

A typical configuration is
%%
\editboxtwo{
 \verb|&PARAM_OCEAN_DYN | & \\
 \verb|EQS_TYPE      = "SLAB",| & ; Ocean model; SLAB or FIXED\_INIT\_TEMP \\
 \verb|TIME_DT       = 3600.D0,| & ; Time interval of ocean dynamics \\
 \verb|TIME_DT_UNIT  = "SEC",| & ; Unit for \verb|TIME_DT| \\
 \verb|/ | & \\
}
When \verb|EQS_TYPE="SLAB"| is selected, the ocean temperature evolves according to the heat budget of the slab ocean.
When \verb|EQS_TYPE="FIXED_INIT_TEMP"| is selected, the initially specified ocean temperature is kept fixed.
For the slab ocean model, the temperature is advanced with a forward Euler method using \nmitem{TIME_DT}. 


\subsubsection{Note}
For the slab ocean model, the ocean temperature is assumed to be vertically uniform.
Therefore, a zeroth-order polynomial would be the natural choice in the vertical direction.
However, the current mesh implementation does not support \verb|PolyOrder_v=0|.
Thus, \nmitem{PolyOrder_v} must be set to one or larger.

Although multiple vertical nodal values are therefore retained redundantly, 
the slab-ocean temperature tendency is applied uniformly to all vertical degrees of freedom, so that the slab-ocean calculation is effectively vertically uniform.

%-------------------------------------------------------------------------------
\subsection{Surface fluxes and atmosphere--ocean exchange}
\label{subsec:ocean_surface_flux}
%-------------------------------------------------------------------------------


For the surface-flux calculation, the ocean component receives the near-surface atmospheric state from the atmosphere through the coupler,
including density, pressure, temperature, wind velocity, specific humidity, and the thickness of the lowest atmospheric layer.
The downward direct shortwave and downward diffuse longwave radiative fluxes are also passed from the atmospheric radiation component to the ocean component.


Using these atmospheric quantities together with the ocean surface temperature, 
the ocean component calculates the surface momentum, sensible-heat, latent-heat, and water-vapor fluxes using surface-flux schemes.
The ocean surface temperature and these surface fluxes are then returned
to the atmospheric component through the coupler and used as the lower boundary condition of the atmosphere.


In the current SCALE-DG, the ocean surface fluxes are calculated using the simple bulk surface-flux scheme described in Section~\ref{subsec:physics_sf_simple}.
The same module and its parameters are shared with the atmospheric surface-flux component.
Therefore, the atmospheric surface-flux component should be enabled
and the simple bulk scheme should be selected as follows:
\editboxtwo{
 \verb|&PARAM_ATMOS | & \\
 \verb|ATMOS_PHY_SF_DO = .true.,| & ; Enable the surface-flux component \\
 \verb|/ | & \\
}
\editboxtwo{
 \verb|&PARAM_ATMOS_PHY_SFC | & \\
 \verb|SFCFLX_TYPE = "SIMPLE",| & ; Use the simple bulk surface-flux scheme \\
 \verb|TIME_DT      = 600.D0,|   & ; Time interval of the atmospheric surface process \\
 \verb|TIME_DT_UNIT = "SEC",|    & ; Unit for \verb|TIME_DT| \\
 \verb|/ | & \\
}
The parameters of the simple bulk surface-flux scheme are specified in \namelist{PARAM_ATMOS_PHY_SF_SIMPLE}. 
See Section~\ref{subsec:physics_sf_simple} for details.
